\usepackage{amsmath,amssymb,amsthm}
\usepackage[capitalise]{cleveref}
\usepackage{hyperref}
\usepackage{csquotes}
\usepackage{placeins}
\usepackage{thmtools}

% Theorem names in babel's convention, as a LyX export writes them.
\providecommand{\theoremname}{Theorem}
\theoremstyle{plain}
\newtheorem{thm}{\protect\theoremname}
\newtheorem{lem}[thm]{Lemma}
\newtheorem{prop}[thm]{Proposition}
\newtheorem{cor}[thm]{Corollary}
\newtheorem{claim}{Claim}
\theoremstyle{definition}
\newtheorem{defn}[thm]{Definition}
\newtheorem{assum}{Assumption}
\newtheorem{example}[thm]{Example}
\theoremstyle{remark}
\newtheorem*{rem}{Remark}

% A reference command of the paper's own: it prints the word and adds the
% label's prefix, so `\thmref{divergence}` points at `thm:divergence`.
\newcommand{\thmref}[1]{Theorem~\ref{thm:#1}}

% A proof under another name: it opens a proof and closes it.
\newenvironment{proofsketch}{\begin{proof}[Sketch]}{\end{proof}}

\newcommand{\N}{\mathbb{N}}
\newcommand{\primes}{\mathcal{P}}
\DeclareMathOperator{\li}{li}
